\chapter[1913 --- Richet]{1913: Charles Richet}
\label{ch:1913_charlesrichet}

\section*{Main Contribution}

\textit{Discovered anaphylaxis---the severe allergic reaction that occurs when sensitized individuals are re-exposed to an allergen---establishing the concept of hypersensitivity.}\cite{nobel-1913,lecture-1913}

\section{Official Citation}

for his discovery of anaphylaxis

\section{Background and Historical Context}

Charles Robert Richet was born on August 26, 1850, in Paris, France, into a distinguished family of physicians and scientists. His father, Alfred Richet, was a prominent surgeon and professor at the University of Paris. The younger Richet studied medicine at the University of Paris, receiving his medical degree in 1876. He then pursued a career in experimental physiology, combining clinical medicine with laboratory research in a manner that was characteristic of the French medical tradition.

Richet's interest in the immune system was sparked by his work on the physiological effects of certain toxins. In the 1880s, he and his colleague Th\'eophile Brouardel investigated the toxic effects of jellyfish stings, studying the venom of the Portuguese man-of-war (\emph{Physalia physalis}) and other marine organisms. This work introduced Richet to the concept of anaphylaxis---the paradoxical and dangerous hypersensitivity reactions that could occur when the body was re-exposed to a substance to which it had been previously sensitized.

The concept of anaphylaxis---from the Greek \emph{ana}, against, and \emph{phylaxis}, protection---was, in the late nineteenth century, a deeply puzzling phenomenon. It was well established that the body could develop immunity to certain toxins and pathogens, and that this immunity was generally protective. But there were puzzling exceptions: some animals, when immunized with small doses of a toxin and then re-exposed to a larger dose, did not show enhanced resistance but instead developed severe, often fatal, reactions. This ``hypersensitivity'' or ``untoward reactivity'' was the opposite of what would be expected if the immune response were purely beneficial.

Several observations had hinted at the existence of hypersensitivity reactions before Richet's work. The French physician Paul Brouillard had described fatal reactions to antitoxin serum injections in 1894, and the German pediatrician Clemens von Pirquet had coined the term ``allergy'' (from the Greek \emph{allos}, other, and \emph{ergon}, work) in 1906 to describe the altered reactivity of patients who had received serum therapy. But none of these observers had systematically investigated the phenomenon or demonstrated it under controlled experimental conditions.

\section{The Discovery}

Richet's discovery of anaphylaxis was made through a series of carefully designed experiments on dogs, conducted between 1902 and 1907 at the physiological laboratory of the University of Paris. The experiments were prompted by a practical question: Why did some patients react violently to antitoxin serum injections, while others tolerated the same injections without ill effect?

The key experiment, published in 1902 in collaboration with Paul Portier, involved the sea anemone \emph{Actinia equina}. Richet and Portier had extracted a toxic substance from the sea anemone---which they called ``actinotoxin''---and were studying its effects on dogs. They observed that dogs that received a small, sublethal dose of the toxin and were then re-exposed to a second dose several weeks later did not show the expected immunity. Instead, they developed a severe, often fatal reaction characterized by violent vomiting, diarrhea, convulsions, collapse, and death. The reaction was dose-dependent and occurred within minutes of the second exposure.

Richet named this phenomenon ``anaphylaxis'' to distinguish it from ``prophylaxis'' (protection), which was the expected outcome of immunization. He demonstrated that anaphylaxis was a general phenomenon that could be produced with a variety of substances, including egg white, horse serum, and bacterial toxins. He showed that the susceptibility to anaphylaxis was specific: animals sensitized with one substance showed anaphylactic reactions only when re-exposed to the same substance, not to different substances. He further demonstrated that anaphylaxis could be transferred from sensitized to unsensitized animals by the transfer of serum---indicating that the ``anaphylactic substance'' was a humoral factor in the blood.

Richet's experimental findings can be summarized as follows:

\textbf{Sensitization.} A small, sublethal dose of a foreign protein (the ``allergen'') administered to an animal produced a state of sensitization that was not associated with any visible symptoms. The animal appeared healthy and showed no signs of immune response.

\textbf{Elicitation.} When the sensitized animal was re-exposed to the same foreign protein, even in small doses, it developed a severe, potentially fatal reaction within minutes. The symptoms included bronchospasm (contraction of the airways), vasodilation and increased vascular permeability (leading to a drop in blood pressure and fluid leakage into tissues), gastrointestinal disturbance, and in severe cases, cardiovascular collapse and death.

\textbf{Specificity.} The anaphylactic reaction was specific: only the substance used for sensitization could elicit the reaction. Sensitization with egg white did not produce anaphylaxis when the animal was re-exposed to horse serum, and vice versa.

\textbf{Transfer by serum.} Anaphylaxis could be passively transferred from a sensitized animal to an unsensitized animal by the transfer of serum. This demonstrated that the ``anaphylactic sensitivity'' was carried by a soluble substance in the blood---a substance that was later identified as immunoglobulin E (IgE), the antibody class responsible for allergic reactions.

Richet's discovery of anaphylaxis was initially met with skepticism and even hostility from some members of the medical establishment. The concept that the immune response could cause disease---rather than prevent it---was deeply counterintuitive to a scientific community that had been trained to think of immunity in exclusively protective terms. Richet himself noted that his paper on anaphylaxis was rejected by several journals before it was finally published. But the experimental evidence was overwhelming, and the reality of anaphylaxis was quickly confirmed by other researchers.

\section{Impact on Modern Medicine}

Richet's discovery of anaphylaxis laid the foundation for the modern understanding of allergic diseases---a field that now affects an estimated 30 percent of adults and 40 percent of children in developed countries. Allergic rhinitis (hay fever), allergic asthma, atopic dermatitis (eczema), food allergies, drug allergies, and insect sting allergies are all manifestations of the hypersensitivity mechanisms that Richet first described.

The understanding of anaphylaxis has led directly to the development of treatments for allergic diseases. Antihistamines, which block the effects of histamine released during allergic reactions, are among the most widely used medications in the world. Epinephrine (adrenaline), which reverses the life-threatening symptoms of anaphylactic shock, is the standard emergency treatment for severe allergic reactions and is carried by millions of people with severe allergies in the form of auto-injectors (such as the EpiPen).

The understanding of the immunological basis of allergy has also led to the development of allergen-specific immunotherapy (desensitization therapy), in which patients are gradually exposed to increasing doses of the allergen that triggers their allergic response. This treatment, which is effective for insect sting allergies, allergic rhinitis, and some forms of asthma, is based on the principle of modifying the immune response to an allergen---a principle that can be traced directly to Richet's observations on anaphylaxis.

More broadly, Richet's work demonstrated that the immune response could be a double-edged sword---protective in some circumstances but harmful in others. This insight has had significant implications for the understanding of autoimmune diseases, in which the immune system attacks the body's own tissues, and for the development of immunosuppressive therapies that modulate the immune response in autoimmune and inflammatory conditions.

\section{Legacy}

Charles Richet received the Nobel Prize in Physiology or Medicine in 1913 ``for his discovery of anaphylaxis.'' He was 63 years old and had been working on the phenomenon for over a decade. The award was a fitting recognition of one of the major discoveries in the history of immunology.

Richet continued his research after receiving the Nobel Prize, publishing extensively on physiology, immunology, and a wide range of other topics. He was a prolific and eclectic scientist whose interests extended beyond medicine to include psychology, philosophy, and literature. He died on December 4, 1935, at the age of 85.

Richet's discovery of anaphylaxis changed the understanding of the immune system from a purely defensive mechanism to a complex system with both beneficial and harmful potential. The allergic diseases that affect hundreds of millions of people worldwide today are all manifestations of the phenomenon that Richet identified, and the treatments for these conditions---from antihistamines to immunotherapy to emergency epinephrine---are all descendants of his pioneering work. In recognizing the destructive potential of the immune response, Richet opened a new chapter in the understanding of human health and disease.


\section{Modern Developments}

Richet's discovery of anaphylaxis has led to the modern allergy field. Epinephrine auto-injectors (EpiPen) are life-saving treatments for anaphylaxis. Understanding of IgE-mediated hypersensitivity has enabled the development of anti-IgE therapy (omalizumab) for severe asthma and food allergies. Oral immunotherapy (e.g., Palforzia for peanut allergy) is the first FDA-approved treatment for food allergy. Dupilumab, targeting IL-4/IL-13, has revolutionized treatment of atopic dermatitis and asthma.


\section{Bibliography}

\begin{thebibliography}{9}
\bibitem{nobel-1913}
Nobel Prize Outreach, ``The Nobel Prize in Physiology or Medicine 1913,''\emph{NobelPrize.org}.\\
\url{https://www.nobelprize.org/prizes/medicine/1913/summary/}

\bibitem{lecture-1913}
Charles Richet, ``Nobel Lecture,''\emph{NobelPrize.org}. Winner-authored primary source; downloadable from the official Nobel Prize archive.\\
\url{https://www.nobelprize.org/prizes/medicine/1913/richet/lecture/}
\end{thebibliography}
